\documentclass[12pt]{article}

%% Basic document formatting
\usepackage{amsmath, amsthm, amssymb, amsfonts}
\usepackage{mathtools}
\usepackage{xspace}
\usepackage{thmtools}
\usepackage{graphicx}
\usepackage{setspace}
\usepackage{fancyhdr}
\usepackage{titling}
\usepackage[left=0.4in,right=0.4in,top=1in,bottom=1in]{geometry}
\usepackage{float}
\usepackage{hyperref}
\usepackage{mathrsfs}
\usepackage{tabularx}
\usepackage[utf8]{inputenc}
\usepackage[english]{babel}
\usepackage{framed}
\usepackage[dvipsnames]{xcolor}
\usepackage{environ}
\usepackage{tcolorbox}
\tcbuselibrary{theorems,skins,breakable}

\usepackage{enumitem}
\setlist[enumerate]{leftmargin=*}
\setlist[enumerate,1]{labelindent=\parindent}
\setlist[enumerate,2]{labelindent=0pt}

% Blackboard Bold
\newcommand{\Z}{\mathbb{Z}}                 % integers
\newcommand{\Zpl}{\mathbb{Z}_{+}}           % positive integers
\newcommand{\Znn}{\mathbb{Z}_{\geq 0}}      % nonnegative integers. 
\newcommand{\Q}{\mathbb{Q}}                 % rationals
\newcommand{\Qpl}{\mathbb{Q}_{+}}           % positive Rationals
\newcommand{\R}{\mathbb{R}}                 % reals
\newcommand{\Rpl}{\mathbb{R}_{+}}           % positive Reals
\newcommand{\C}{\mathbb{C}}                 % complex numbers
\newcommand{\F}{\mathbb{F}}                 % field

% Words
\newcommand{\st}{\text{ such that }}
\newcommand{\wrt}{\text{ with respect to }}
\newcommand{\with}{\text{ with }}
\newcommand{\ie}{\text{, i.e. }}
\newcommand*{\ditto}{\text{---}\texttt{"}\text{---}}

% Operators
\newcommand{\abs}[1]{\left|#1\right|}                   % absolute value
\newcommand{\norm}[1]{\abs{\abs{#1}}}
\newcommand{\floor}[1]{\left\lfloor #1 \right\rfloor}   % floor
\newcommand{\ceil}[1]{\left\lceil #1 \right\rceil}      % ceiling
\newcommand{\powset}[1]{\mathscr{P}(#1)}

% Category Theory 
\renewcommand{\hom}{\text{Hom}}
\newcommand{\homspace}[3]{\hom_{#1}(#2, #3)}
\newcommand{\coproduct}[9]{
  \begin{tikzcd}[ampersand replacement=\&]
      \& \& #1 \arrow[dll, bend right, "#6"] \arrow[dl, "#5"] \\
   #4 \& #3 \arrow[l, "#9"] \\
      \& \& #2 \arrow[ull, bend left, "#8"] \arrow[ul, "#7"]
  \end{tikzcd}
}

\newcommand{\product}[9]{ 
  \begin{tikzcd}[ampersand replacement=\&]
      \& \& #3 \\
      #1 \arrow[r, "#7"] \arrow[urr, bend left, "#5"] \arrow[drr, bend right, "#6"] \& #2 \arrow[ur, "#8"] \arrow[dr, "#9"] \\ 
      \& \& #4
  \end{tikzcd}
}


% Algebra 
\newcommand{\Syl}[2]{\text{Syl}_{#1}{(#2)}}           % set of Sylow-p subgroups 
\newcommand{\<}{\langle}                % \<x\>, subgroup generated by x 
\renewcommand{\>}{\rangle}
\renewcommand{\index}[2]{[#1 : #2]}
\newcommand{\id}{\text{id}}             % identity element
\newcommand{\lhdneq}{%
  \mathrel{\ooalign{$\lneq$\cr\raise.22ex\hbox{$\lhd$}\cr}}}
\newcommand{\order}[1]{\text{o}(#1)}    % order of an element
\newcommand{\spec}{\operatorname{Spec}}
\newcommand{\rad}{\operatorname{rad}}   % radical of an ideal 
\newcommand{\sym}[1]{\mathfrak{S}_{#1}}
\newcommand{\aut}[1]{\text{Aut}(#1)} 
\newcommand{\gal}[2]{\text{Gal}(#1 / #2)} 
\newcommand{\inn}[1]{\text{Inn}(#1)} 
\let\oldcong\cong
\let\oldequiv\equiv
\renewcommand{\cong}{\oldequiv}
\renewcommand{\equiv}{\oldcong}

\newcommand{\triangleleftneq}{\mathrel{\ooalign{%
  $\trianglelefteq$\cr
  \hidewidth\raise0.06ex\hbox{$\not\mkern4mu$}\hidewidth\cr
}}}

\makeatletter % cycle
\newcommand{\cyc}[1]{(\mathbf{\cyc@process#1\relax})}
\def\cyc@process#1#2\relax{%
	#1%
	\ifx\relax#2\relax
	\else
		\,\cyc@process#2\relax
	\fi
}
\makeatother

\newcommand{\GL}[2]{\text{GL}_{#1}(#2)}                             % general linear group
\newcommand{\SL}[2]{\text{SL}_{#1}(#2)}                             % special linear group
\newcommand{\gl}[2]{\mathfrak{gl}_{#1}(#2)}
\renewcommand{\sl}[2]{\mathfrak{sl}_{#1}(#2)}
\newcommand{\so}[2]{\mathfrak{so}_{#1}(#2)}
\newcommand{\su}[1]{\mathfrak{su}(#1)}

% Linear Algebra
\newcommand{\bmat}[1]{\begin{bmatrix}#1\end{bmatrix}}   % bracketed matrix
\newcommand{\rank}{\operatorname{rank}}                 % rank
\newcommand{\nullity}{\operatorname{nullity}}           % nullity
\newcommand{\tr}{\operatorname{tr}}

% Combinatorics
\newcommand{\qnum}[1]{\left[#1\right]_q}                % q-analog
\newcommand{\qfact}[1]{\left[#1\right]!_q}              % q-analog factorial 
\newcommand{\qbinom}[2]{\binom{#1}{#2}_q}               % Gaussian binomial coefficient

% Topology/Analysis
\newcommand{\ball}[2]{\text{B}_{#1}(#2)}  % B_r(x): open r-balls around x
\newcommand{\diam}{\text{diam}}           % diamter of a set in metric space 
\newcommand{\lowint}[2]{\underline{\int_{#1}^{#2}}}
\newcommand{\upint}[2]{\overline{\int}_{#1}^{#2}}

% Misc Notation
\newcommand{\oneton}[1]{\{1, \ldots, #1\}}
\newcommand{\defeq}{\vcentcolon=}   % :=
\newcommand{\eqdef}{=\vcentcolon}   % =: 
\renewcommand{\bf}[1]{\textbf{#1}}
\newcommand{\im}{\operatorname{im}}
\newcommand{\fnrest}[2]{\left.#1\right|_{#2}}
\newcommand{\isom}{\overset{\sim}{\rightarrow}} 


\renewcommand{\thesection}{\arabic{section}.}
\renewcommand{\thesubsection}{(\alph{subsection})}

\newtheorem{claim}{Claim}
\newtheorem{theorem}{Theorem}
\newtheorem{proposition}{Proposition}
\newtheorem*{lemma}{Lemma}

%% Headers & title setup
\newcommand{\course}{Math140B}
\newcommand{\myname}{Jay Ser}
\setlength{\headheight}{14.5pt}
\pagestyle{fancy}
\fancyhf{}
\renewcommand{\headrulewidth}{0.4pt}
\lhead{\course}
\rhead{\myname}
\cfoot{\thepage}
\setlength{\droptitle}{-4em} 
\title{\course\ - HW \#10}
\author{\myname}
\date{2026.03.15} 

\begin{document}
\maketitle
\thispagestyle{fancy}

%------------------------------------------------------------------------------%

\section*{Q12} 
\subsection{}
If $m \neq 0$, then 
\begin{align*}
  c_m & = \frac{1}{2\pi} \int_{-\pi}^{\pi} f(x) e^{-imx} dx \\ 
      & = \frac{1}{2\pi} \int_{-\delta}^{\delta} e^{-imx} \\ 
      & = -\frac{1}{2m\pi i} \left.e^{-imx}\right|_{-\delta}^{\delta} \\ 
      & = \frac{1}{2m\pi i} (e^{im\delta} - e^{-im\delta}) \\ 
      & = \frac{\sin(m\delta)}{m\pi}
\end{align*}
For the special case $m = 0$, 
\begin{align*}
  c_0 & = \frac{1}{2\pi} \int_{-\pi}^{\pi} f(x) e dx \\ 
      & = \frac{1}{2\pi} \int_{-\delta}^{\delta} dx \\ 
      & = \frac{\delta}{\pi}
\end{align*}

\subsection{}
Since $f$ is differentiable at $x_0 = 0$, Theorem 8.14 applies, yielding 
\begin{equation} \label{eq:12_b_thm}
  \lim_{N \rightarrow \infty} \sum_{n = -N}^{N} c_n e^{in x_0} = \lim_{N \rightarrow \infty} \sum_{n = -N}^{N} c_n = f(x_0) = 1
\end{equation}
Namely, $\sum_{n = 1}^{\infty} (c_n + c_{-n})$ converges. 
Rearranging \eqref{eq:12_b_thm}and using the values for the $c_m$'s calculated in (a), 
\begin{align*} 
  1 - \frac{\delta}{\pi} & = 1 - c_0 \\ 
                         & = \sum_{n = 1}^{\infty} (c_n + c_{-n}) \\ 
                         & = \sum_{n = 1}^{\infty} (\frac{1}{n\pi} \sin(n\delta) + \frac{1}{-n\pi} \sin(-n\delta)) \\ 
                         & = \frac{2}{\pi} \sum_{n = 1}^{\infty} \frac{\sin(n\delta)}{n},
\end{align*}
which is equivalent to showing 
$$\sum_{n = 1}^{\infty} \frac{\sin(n\pi)}{n} = \frac{\pi - \delta}{2}$$

\subsection{}
$f$ has period $2\pi$, so Parseval's Theorem applies, giving 
$$\frac{1}{2\pi} \int_{-\pi}^{\pi} \abs{f(x)}^2 dx = \sum_{n = -\infty}^{\infty} \abs{c_n}^2$$
Simplifying,
\begin{align*}
  \frac{\delta}{\pi} & = \frac{\delta^2}{\pi^2} + \sum_{n = 1}^{\infty} \left( \frac{(\sin(n\delta))^2}{n^2\pi^2} + \frac{(\sin(-n\delta))^2}{n^2 \pi^2} \right) \\ 
  \frac{\delta}{\pi}(1 - \frac{\delta}{\pi}) & = \frac{2}{\pi^2} \sum_{n = 1}^{\infty} \frac{(\sin(n\delta))^2}{n^2} \\ 
  \frac{\pi - \delta}{2} & = \sum_{n = 1}^{\infty} \frac{(\sin(n\delta))^2}{n^2\delta},
\end{align*}
as was to be shown. 

\section*{Q13}
Let $f(x) = x$ for $0 \leq x < 2\pi$ with a period of $2\pi$. 
When $n = 0$, 
$$c_0 = \frac{1}{2\pi} \int_{0}^{2\pi} x dx = \frac{1}{2\pi} \left. \frac{x^2}{2}\right|_{0}^{2\pi}  = \pi$$
For $n \neq 0$, take $H(x) = x$, $G(x) = \frac{1}{-in} e^{-inx}$ and perform integration by parts to obtain 
\begin{align*}
  c_n & = \frac{1}{2\pi} \int_{0}^{2\pi} x e^{-inx} dx \\ 
      & = \frac{1}{2\pi} (\frac{2\pi}{-in} + 0 + \frac{-1}{n^2} \left.e^{inx}\right|_{0}^{2\pi}) \\ 
      & = \frac{i}{n}
\end{align*}
Parseval's Theorem states 
\begin{equation} \label{eq:13_parseval}
  \frac{1}{2\pi} \int_{0}^{2\pi} \abs{f(x)}^2 dx \underset{\triangle}{=} \sum_{n = -\infty}^{\infty} \abs{c_n}^2 = \pi^2 + \sum_{n = 1}^{\infty} \frac{2}{n^2}
\end{equation}
Evaluating the integral on the left hand side, 
\begin{align*}
  \frac{1}{2\pi} \int_{0}^{2\pi} \abs{f(x)}^2  
  & = \frac{1}{2]pi} \int{0}^{2\pi} x^2 dx \\ 
  & = \frac{1}{2\pi} \left.\frac{x^3}{3}\right|_{0}^{2\pi} \\ 
  & = \frac{4\pi^2}{3}
\end{align*}
Thus rearranging \eqref{eq:13_parseval} yields 
$$\sum_{n = 1}^{\infty} \frac{1}{n^2} = \frac{\pi^2}{6}$$

\section*{Q14}
Let $f(x) = (\pi - \abs{x})^2$ on $[-\pi, \pi]$. 
Since $f(-\pi) = f(\pi)$, extend the domain of $f$ to $\R$ such that $f$ has a period of $2\pi$. 
When $n = 0$, 
\begin{align*}
  a_0 & = \frac{1}{2\pi} \int_{-\pi}^{\pi} (\pi - \abs{x})^2 dx \\ 
      & = \frac{1}{2\pi} \left( \int_{0}^{\pi} (\pi^2 - 2\pi x + x^2) dx + \int_{-\pi}^{0} (\pi^2 + 2\pi x + x^2)dx \right)\\ 
      & = \frac{1}{2\pi} \left( \left[ \pi^2 x - \pi x^2 + \frac{1}{3}x^3 \right]_{0}^{\pi} + \left[ \pi^2 x + \pi x^2 + \frac{1}{3}x^3 \right]_{-\pi}^{0} \right)\\ 
      & = \frac{\pi^2}{3}
\end{align*}
When $n > 0$, $b_n = \int_{-\pi}^{\pi} f(x) \sin(nx)dx = 0$ since $f(x)\sin(nx)$ is an odd function. 
Next, using the fact that $f(x) \cos(nx)$ is even, one can calculate  
\begin{align*}
  a_n & = \frac{1}{\pi} \int_{-\pi}^{\pi} (\pi - \abs{x})^2 \cos(nx) dx \\ 
      & = \frac{2}{\pi} \int_{0}^{\pi} (\pi^2 - 2\pi x + x^2)\cos(nx)dx \\  
      & = \frac{2}{\pi} \left[ \frac{n^2 (\pi - x)^2 - 2)\sin(nx) - 2n(\pi - x)\cos(nx)}{n^3} \right]_{0}^{\pi}  \\ 
      & = \frac{4}{n^3 \pi} (\pi n - \sin(n\pi)) \\ 
      & = \frac{4}{n^2}
\end{align*}
So the Fourier series of $f$ is 
$$f(x) \sim \frac{\pi^2}{3} + \sum_{n = 1}^{\infty} \frac{4}{n^2} \cos(nx)$$
For $x_0 \neq 0$, $f$ is differentiable, so Theorem 8.14 applies. 
When $x_0$, it's clear that for some small $t$, there is $M$ big enough to satisfy $\abs{f(t) - f(0)} \leq M\abs{t}$, for example $M > 2\pi + \abs{t}$, so Theorem 8.14 applies everywhere, yielding 
\begin{equation} \label{eq:14_res}
  f(x) = \frac{\pi^2}{3} + \sum_{n = 1}^{\infty} \frac{4}{n^2}\cos(nx)
\end{equation}
on all of $\R$. 
Applying \eqref{eq:14_res} to $x = 0$, $\pi^2 = \frac{\pi^2}{3} + \sum_{n = 1}^{\infty} \frac{4}{n^2}$, i.e., 
$$\frac{\pi^2}{6} = \sum_{n = 1}^{\infty} \frac{1}{n^2}$$

To get the second identity, apply Parseval's theorem since $f$ is only discontinuous at $x = 0$ on $[-\pi, \pi]$ and hence $f \in \mathscr{R}_{-\pi}^{\pi}$: 
\begin{equation} \label{eq:14_2nd_parseval}
  \frac{1}{2\pi} \int_{-\pi}^{\pi} \abs{f(x)}^2 dx = \sum_{n = -\infty}^{\infty} \abs{c_n}^2
\end{equation}
The right hand side simplifies to 
$$\abs{a_0}^2 + 2 \sum_{n = 1}^{\infty} \abs{\frac{a_n}{2}}^2 = \frac{\pi^4}{9} + 8 \sum_{n = 1}^{\infty} \frac{1}{n^4}$$
Meanwhile, the left hand side is 
\begin{align*}
  \frac{1}{2\pi} \left[ \int_{0}^{\pi} (\pi - x)^4 dx + \int_{-\pi}^{0} (\pi + x)^4 dx \right] 
  & = \frac{1}{2\pi} \left[ \int_{-\pi}^{0} y^4 dx + \int_{0}^{\pi} y^4 dx \right] \\ 
  & = \frac{1}{2\pi} \left( \frac{1}{5} \left.y^5\right|_{-\pi}^{0} + \frac{1}{5}\left.y^5\right|_{0}^{\pi} \right) \\ 
  & = \frac{\pi^4}{5} 
\end{align*}
So \eqref{eq:14_2nd_parseval} is equivalent to $\frac{\pi^4}{5} - \frac{\pi^4}{9} = 8\sum_{n = 1}^{\infty} \frac{1}{n^4}$,
i.e., 
$$\frac{\pi^4}{90} = \sum_{n = 1}^{\infty} \frac{1}{n^4}$$
\end{document}
